\centerline{\bf \Huge Тренировочная олимпиада}
\centerline{\bf \Huge Муницип 9 класс}

\begin{flushright}
        \scriptsize{Каждая задача оценивается в 4 балла.}
\end{flushright}

\problem{9.1} У каждого двузначного числа нашли сумму цифр, а потом все эти суммы перемножили. На сколько нулей заканчивается полученное произведение?

\problem{9.2} Можно ли расставить по кругу 8 различных натуральных чисел так, чтобы для любых двух соседей отношение большего к меньшему было целым, и все такие отношения были различны и не превосходили 10 ?

\problem{9.3} Дано 6 отрезков разной длины таких, что из любых трёх можно составить треугольник. Докажите, что из всех шести отрезков можно составить 2 непрямоугольных треугольника, использовав каждый отрезок ровно один раз.

\problem{9.4} Числа $a$ и $b$ таковы, что уравнение $x^2+a x+b=0$ имеет корни $x_1$ и $x_2, x_1<x_2$. Известно, что, если заменить в этом уравнении коэффициент на любое число из отрезка $\left[x_1, x_2\right]$, то полученное уравнение всё равно будет иметь корни. Докажите, что тогда оно будет иметь корни при всех $a$.

\problem{9.5} Серединные перпендикуляры к сторонам $A C$ и $B C$ неравнобедренного остроугольного треугольника $A B C$ пересекают высоту из точки $C$ или её продолжение в точках $X$ и $Y$. Пусть $O$ --- центр описанной окружности $\triangle A B C$. Докажите, что $C O$ касается описанной окружности $\triangle O X Y$.